\chapter{Fundamentos teóricos}\label{ch:theoretical_background}

% SECTION 1

\section{Ciclones extratropicales en el Atlántico Sur}\label{sec:tb_fundamentos}
La teoría clásica (e.g., modelo noruego) define a los ciclones extratropicales como vórtices de núcleo frío, y climatologías robustas como la de \citet{hoskins2005new} demuestran que su génesis global está dominada por la inestabilidad baroclínica de latitudes medias, la caracterización de estos sistemas en la región específica del Atlántico Sur requiere matices dinámicos adicionales. \citet{sinclair1995climatology} consolidó esta perspectiva al describir las características del ciclo de vida de los ciclones en el Hemisferio Sur, resaltando que la ciclogénesis está influenciada por los gradientes térmicos oceánicos y la orografía. De manera específica, \citet{inatsu2004zonal} demostraron que la asimetría zonal del \textit{storm track} en el Hemisferio Sur está forzada por la topografía localizada de los Andes. Este forzamiento orográfico modula la propagación de ondas baroclínicas sobre el subcontinente sudamericano \citep{vera2002cold}, definiendo las regiones preferenciales de desarrollo en latitudes subtropicales y extratropicales. Esta variabilidad estructural y su dependencia de zonas con baroclinidad han sido analizadas por \citet{simmonds1999southern} aunque para latitudes altas.

En este contexto, la transición de una perspectiva euleriana a una lagrangiana resulta necesaria; \citet{sinclair1994objective} demostró que el uso de la vorticidad relativa a 850 hPa es físicamente superior a la presión a nivel del mar para identificar sistemas en etapas tempranas, ya que permite aislar la rotación de mesoescala frente al flujo zonal del Hemisferio Sur. En su tesis doctoral, \citet{coutodesouza2024thesis} realiza una revisión de la física de estos sistemas, concluyendo que la definición adiabática estándar es insuficiente para la región. Si bien la inestabilidad baroclínica proporciona el entorno favorable para la génesis, la evolución y la intensidad del sistema puede ser modulada por procesos diabáticos asociados a la liberación de calor latente. Este comportamiento es ilustrado por los experimentos de sensibilidad numérica presentados por \citet{reboita2022from}, donde la supresión de los flujos de calor latente altera la estructura térmica y la intensidad del vórtice.

% SECTION 2

\section{Regiones de ciclogénesis}\label{sec:tb_regiones}

A partir del análisis de la vorticidad relativa en niveles bajos, \citet{reboita2010south} establecieron la zonificación climática de estos procesos, delimitando las regiones de la costa de Uruguay y el sur de Brasil, así como la región costera de Argentina, como los principales focos de actividad. Mientras que los sistemas extratropicales puros dominan las latitudes altas, \citet{gozzo2014subtropical} demuestran que en las regiones de latitudes medias y subtropicales (20°S--35°S) existe una alta frecuencia de génesis de ciclones híbridos. La interacción entre el flujo de los oestes, la topografía de los Andes y la Temperatura Superficial del Mar (TSM) define tres regiones de génesis con características físicas diferenciadas \citep{gramcianinov2019properties}: 

\subsection{Sur-sureste de Brasil (SBR)}\label{subsec:tb_sbr}

Según \citet{reboita2022from}, la región de génesis RG1 —cuya equivalencia con la SBR fue establecida en la Sección \ref{sec:definicion_regiones}— constituye una de las regiones preferenciales de ciclogénesis en la costa sur/sudeste de Brasil, y el caso de estudio del ciclón Raoni ilustra cómo los sistemas que actúan en este sector pueden presentar características subtropicales durante su ciclo de vida. Estos sistemas son llamados como ciclones "híbridos", porque aunque se inician por inestabilidad baroclínica, su intensificación está modulada por procesos diabáticos y la inestabilidad de la capa límite marina, lo que favorece el desarrollo de convección en el sector cálido \citep{marrafon2022classificacao}. 
El análisis energético de \citet{coutodesouza2024thesis} revela que, a diferencia de las regiones australes, los ciclones en SBR presentan una conversión de energía "mixta", donde los procesos diabáticos (liberación de calor latente) juegan un rol tan importante como la baroclinicidad. Por otro lado, \citet{gramcianinov2024early} demuestran que, durante la ciclogénesis temprana en la SBR, los modelos de alta resolución logran representar las estructuras de mesoescala de advección cálida y convergencia de humedad en capas bajas, a diferencia de los modelos de baja resolución, que sistemáticamente suavizan estos mecanismos. Este resultado valida que la dinámica híbrida de SBR no es un artefacto estadístico sino una señal física.

\subsection{Cuenca del Río de la Plata (LPB)}\label{subsec:tb_lpb}

La región comprendida entre $30^\circ$S y $40^\circ$S constituye uno de los principales \textit{hotspots} de ciclogénesis en el Atlántico Sudoccidental. Su mecanismo dominante es la ciclogénesis de sotavento: como establecen \citet{gan1994influence}, el flujo de los oestes experimenta una compresión y posterior estiramiento vertical de la columna de atmosfera al atravesar los Andes. Por conservación de la vorticidad potencial, este estiramiento en el flanco oriental induce una caída de presión en superficie que favorece la iniciación de la rotación ciclónica. 

En este contexto, \citet{reboita2010south} muestran que esta región se caracteriza por una alta frecuencia de sistemas ciclónicos y una interacción con el transporte de humedad desde latitudes tropicales, mientras que el aporte de aire cálido y húmedo asociado al Chorro de Capas Bajas (SALLJ) resulta determinante para la intensidad de la precipitación \citep{crespo2020potential}.

Consistentemente, \citet{gramcianinov2024early} confirman mediante modelos regionales que la convergencia de humedad en capas bajas constituye el mecanismo dominante de intensificación en LPB, proyectando además que bajo escenarios de calentamiento global este forzante se reforzará hacia fines de siglo. Si bien la frecuencia total de ciclones en la región disminuirá, los sistemas que persistan exhibirán flujos de humedad y advección de calor más intensos en sus etapas tempranas, lo que sugiere una mayor severidad potencial de los eventos intensos de precipitación asociados.

\subsection{Argentina (ARG)}\label{subsec:tb_arg}

El sector austral ($>40^\circ$S, o su similar RG3) responde a un régimen dinámico clásico dominado por la baroclinicidad de gran escala. A diferencia de los focos subtropicales, \citet{hoskins2005new} demuestran que el desarrollo en las latitudes australes está estrictamente acoplado en toda la columna troposférica, siendo gobernado por la inestabilidad baroclínica asociada al chorro polar y la propagación de ondas de Rossby a lo largo del storm track de alta latitud. Aquí, \citet{simmonds2000variability} indican que los sistemas ciclónicos dependen menos de forzantes locales de superficie y más de la dinámica de la alta troposfera. Como resultado, estos sistemas tienden a presentar una evolución más rápida y directamente acoplada al flujo de niveles altos, en contraste con las regiones subtropicales del sudeste de Brasil \citep{vera2002cold}. 

% SECTION 3

\section{Ciclo de vida y Modelo conceptual}\label{sec:tb_evolucion}

\subsection{Modelo conceptual}\label{subsec:tb_modelos}

La distribución espacial de las variables en superficie (viento y precipitación) está condicionada por el modelo conceptual que describe la evolución del ciclón. Frente al modelo clásico los ciclones extratropicales intensos del Atlántico Sur suelen seguir el modelo de Fractura Frontal propuesto por \citet{shapiro1990life}. En el Hemisferio Norte la revisión histórica de \citet{schultz2021antecedents}, aunque el modelo Noruego precedió al de Shapiro-Keyser por siete décadas, los antecedentes de la fractura frontal ya eran observables en la literatura de la Escuela de Bergen, siendo el modelo de Shapiro-Keyser una síntesis de elementos estructurales que esa escuela había identificado pero no formalizado. Este modelo se distingue por cuatro rasgos diagnósticos: la fractura del frente frío (\textit{frontal fracture}), el frente curvado (\textit{bent-back front}), la estructura en T (\textit{T-bone}) y el aislamiento de una masa de aire cálido en el núcleo (\textit{warm seclusion}).

Para comprender la dinámica interna que rige la asimetría de estas variables, resulta fundamental complementar este paradigma con el modelo de cinturones de transporte \citep{browning1986conceptual}. En su adaptación conceptual para el Hemisferio Sur, Silva Dias y Hallak (2012) detallan cómo el WCB y su respectiva corriente en chorro de bajo nivel se organizan en torno al vórtice. El ascenso y giro anticiclónico de este flujo cálido modula las estructuras de mesoescala, concentrando las bandas de precipitación y el cizallamiento del viento por delante del frente frío. 
Desde una perspectiva global, \citet{catto2015fronts} demuestran que hasta el 69\% de los WCBs están asociados con frentes identificables objetivamente y los frentes asociados a WCBs son entre 2 y 10 veces más propensos a producir precipitación extrema que los frentes sin WCB. 
Esta vinculación directa entre la dinámica frontal y la precipitación extrema, y no meramente la intensidad del vórtice central, justifica el enfoque espacial de esta tesis.

\subsection{Fases del ciclo de vida}\label{subsec:tb_fases}
Para caracterizar la evolución temporal del viento y la precipitación, este estudio adopta la segmentación del ciclo de vida en cuatro fases discretas. Esta clasificación sigue la metodología objetiva del paquete computacional \textit{CycloPhaser} \citep{deSouza2025CycloPhaser}, aplicado por \citet{coutodesouza2024new}, el cual utiliza la tasa de cambio de la vorticidad relativa para definir los siguientes regímenes estructurales. Cabe notar que cada fase no es homogénea internamente: los instantes próximos a sus transiciones comparten rasgos de la etapa adyacente, por lo que la condición estructural más representativa de cada régimen dinámico corresponde al núcleo temporal de la fase, alejado de ambos bordes de transición.

\begin{enumerate}
    \item Incipiente: Corresponde a la aparición inicial del núcleo de vorticidad organizada. En el contexto regional, esta fase está frecuentemente asociada a la interacción entre una vaguada en altura y forzantes de superficie, como la orografía de los Andes o gradientes térmicos costeros.    
    
    \item Intensificación: Período caracterizado por el rápido crecimiento de la vorticidad ciclónica. Según el análisis energético de \citet{coutodesouza2024thesis}, esta etapa es impulsada por una retroalimentación entre la inestabilidad dinámica y los flujos de calor latente y sensible.
    
    \item Madurez: Etapa donde el sistema alcanza su intensidad máxima (pico de vorticidad) y su mayor extensión horizontal. Es durante esta fase donde los ciclones oceánicos intensos consolidan la evolución estructural documentada observacionalmente por \citet{shapiro1999bridge}. En esta configuración, la deformación del flujo genera la fractura del frente frío y envuelve una masa de aire para que quede aislada en el núcleo del vórtice (\textit{warm seclusion}). Al completarse este proceso las huellas de precipitación y viento quedan desacopladas.
    
    \item Decaimiento: Fase final marcada por el debilitamiento progresivo del vórtice. Este proceso ocurre debido a la fricción superficial y al cese de la conversión de energía baroclínica una vez que el sistema se ha ocluido completamente o ha entrado en una zona barotrópica.
\end{enumerate}


\subsection{Estructuras de mesoescala generadoras de vientos extremos}\label{subsec:tb_mesoescala}

Dentro del paradigma de Shapiro-Keyser, una estructura de mesoescala de especial relevancia para los vientos extremos en superficie es el Sting Jet (SJ). Mediante simulaciones de alta resolución, \citet{clark2005sting} caracteriza su estructura tridimensional completa como un flujo coherente que desciende desde niveles medios, acelera de menos de 20~m~s$^{-1}$ a más de 45~m~s$^{-1}$ en unas 4 horas, y produce la región de vientos superficiales más dañinos en la zona de fractura frontal del ciclón —distinta del WCB y del CCB—. Posteriormente, la revisión sistemática de \citet{clark2018sting} consolida esta definición, precisando que el SJ es un flujo descendente de niveles medios que actúa durante pocas horas en ciclones tipo Shapiro-Keyser, generando una región distintiva de vientos en superficie en la misma zona identificada por \citet{clark2005sting}.

La distinción observacional entre el CCB y el SJ fue establecida por \citet{martinez2014cold}, quienes mediante trazadores químicos demostraron que ambas corrientes constituyen masas de aire físicamente distintas, el CCB se identifica como una corriente de larga duración que acelera en el flanco frío del frente curvado. El SJ, en cambio, se origina en la cabeza nubosa y experimenta una aceleración rápida durante su descenso hacia la región de fractura frontal, generando vientos intensos y localizados.

Aunque la climatología global de eventos tipo SJ permanece incompleta, siendo identificada por \citet{clark2018sting} como una brecha en el conocimiento actual, su relevancia para el Atlántico Sur radica en que esta cuenca presenta una alta frecuencia de ciclones que siguen la evolución de Shapiro-Keyser \citep{gozzo2013air}, constituyendo así el entorno dinámico favorable para su desarrollo. Por otro lado, \citet{hyeok2024extrem} identifican un mecanismo adicional de generación de vientos extremos superficiales en ciclones de latitudes medias: los vientos horizontales que impactan sobre la superficie frontal fría son bloqueados por esta, generando corrientes descendentes (\textit{downdrafts}) que transportan momento horizontal de niveles altos hacia la superficie, intensificando los vientos en los sectores sur y oeste del centro ciclónico (norte y centro para Hemisferio Sur). Este mecanismo —distinto del SJ pero igualmente localizado en la región de fractura frontal— representa una razón más de asimetría en la distribución superficial de los vientos máximos. 

La consideración explícita de estas estructuras de mesoescala provee el marco conceptual para interpretar por qué los extremos de viento en superficie pueden presentar distribuciones espaciales más complejas que las que predice el modelo de escala sinóptica, y justifica metodológicamente el uso de la funcion de densidad de probabilidad (PDF) usando la estimación de densidad de kernel (KDE) bidimensional para mapear la huella real de estos eventos.

% SECTION 4

\section{Vorticidad relativa}\label{sec:tb_vorticidad}

Debido a la extensión oceánica del Hemisferio Sur, los ciclones extratropicales se desarrollan inmersos en un flujo de los oestes. Esta corriente de fondo hace que los sistemas se trasladen con gran rapidez y que, en sus etapas iniciales, la caída de presión se vea enmascarada, presentándose a menudo como vaguadas abiertas en lugar de centros cerrados \citep{sinclair1994objective}. Dada esta rapidez de traslación y complejidad térmica, la presión al nivel del mar (MSLP) resulta una métrica inadecuada para la detección temprana. 

Para resolver este problema y aislar la circulación del sistema, la recomendación metodológica de \citet{sinclair1997objective} es utilizar la vorticidad relativa en 850 hPa ($\zeta_{850}$) como variable de seguimiento. Conceptualmente, la vorticidad relativa mide la rotación macroscópica local del fluido atmosférico respecto a un marco de referencia en la superficie terrestre. Físicamente se define como:

\begin{equation}\label{eq:vorticidad_relativa}
\zeta = \mathbf{k} \cdot (\nabla \times \mathbf{V}) = \frac{\partial v}{\partial x} - \frac{\partial u}{\partial y},
\end{equation}
donde $u$ y $v$ son las componentes zonal y meridional del viento horizontal, respectivamente.

Según \citet{hoskins2005new}, el uso de $\zeta$ permite filtrar el flujo de gran escala y capturar la rotación inducida por estos forzantes de mesoescala antes de que se manifiesten en el campo de presión, permitiendo una identificación más precisa de las fases incipientes, ademas la literatura reciente, como en \citet{gramcianinov2019properties} utiliza $\zeta_{850}$ para identificar con precisión las regiones de ciclogénesis. Asimismo, \citet{padilhareinke2024objective} validan que los algoritmos basados en este nivel isobárico son más eficientes para capturar la posición del ciclón durante su etapa de mayor intensidad. 

Recientemente, el uso de $\zeta_{850}$ se ha extendido más allá del simple seguimiento (tracking) hacia la caracterización estructural. \citet{coutodesouza2024new} demuestran que la evolución temporal de esta variable es el predictor más confiable para segmentar objetivamente el ciclo de vida en fases dinámicas (incipiente, intensificación, madurez y decaimiento). En un contexto global el análisis de clústeres de \citet{corner2025classification} ratifican a la vorticidad relativa en 850 hPa como una de las métricas irreductibles necesarias para describir la intensidad y la estructura de los ciclones extratropicales en el hemisferio norte.

El uso de umbrales percentílicos sobre la distribución de $\zeta_{850}$ para estratificar ciclones por intensidad dinámica cuenta con respaldo metodológico sólido. \citet{priestley2022improved} aplican explícitemente este criterio: seleccionan el 10\% de ciclones con mayor vorticidad ciclónica en su pico de intensidad para construir compositos representativos de la estructura ciclónica. En esta tesis, denominaremos p90 a ese subconjunto (el percentil 90 de vorticidad absoluta), ya que aísla los sistemas más intensos desde el punto de vista dinámico. \citet{andrade2024composite} documentan de forma complementaria, para el Atlántico Sudoccidental, que los sistemas ubicados en el percentil 10 de presión mínima —es decir, los ciclones con presión más baja— exhiben patrones sinópticos diferenciados y una mayor coherencia entre los extremos de viento en superficie y la intensidad del vórtice. Si bien la presión y la vorticidad miden aspectos distintos de la intensidad ciclónica, ambos umbrales percentílicos (p90 para vorticidad y percentil 10 para presión) responden a la misma lógica: retener solo los casos más intensos, lo que permite validar estadísticamente estas submuestras como poblaciones física y estructuralmente homogéneas.

% SECTION 5

\section{Estructura superficial y extremos}\label{sec:tb_compuestos}

La complejidad inherente a los ciclones extratropicales del Atlántico Sur requiere superar el análisis univariado tradicional. En esta sección se fundamenta la adopción del paradigma de \textit{evento compuesto} y se describen las herramientas estadísticas avanzadas (MEOF y PDF) necesarias para capturar su estructura asimétrica y la distribución de sus máximos.

\subsection{Asimetría espacial de los máximos}\label{subsec:tb_kde_asimetria}

Los ciclones extratropicales son sistemas altamente asimétricos, donde los máximos raramente coinciden con el centro de vorticidad. Es importante notar que gran parte de los modelos conceptuales clásicos fueron desarrollados para el Hemisferio Norte; por lo tanto, al analizar el Atlántico Sur, la distribución espacial debe entenderse como una imagen invertida latitudinalmente tomando al ecuador como eje.

\begin{itemize}
    \item Precipitación: Está fundamentalmente controlada por el WCB. En la literatura clásica del Hemisferio Norte, este flujo se describe ascendiendo típicamente hacia el noreste \citep{blanchard2021warm}. Sin embargo, al trasladar este concepto al Hemisferio Sur, el efecto del espejo hace que el WCB concentre la humedad en el sector sureste del sistema durante su etapa de desarrollo. De acuerdo con el modelo conceptual adaptado para el Hemisferio Sur por Silva Dias y Hallak (2012) a partir de \citet{browning1986conceptual}, el ascenso inclinado de este flujo cálido organiza la precipitación en estructuras de mesoescala bien definidas, generando bandas anchas de lluvia estratiforme por delante del sistema y bandas estrechas de convección vigorosa en la interfaz del frente frío. Además, como demuestran \citet{oertel2021warm}, el WCB no solo produce esta precipitación a gran escala, sino que frecuentemente alberga núcleos de convección intensa embebida. La marcada asimetría de este campo es corroborada empíricamente por \citet{naud2020evaluation}, quienes demuestran que la precipitación en los ciclones oceánicos se concentra abrumadoramente en el sector oriental (ascendente y húmedo), mientras que el flanco occidental permanece subsidente y seco.

    \item Viento: Al igual que la precipitación, el campo de viento en superficie presenta una asimetría que cambia con el ciclo de vida del sistema. El mecanismo dinámico central que genera el jet de bajo nivel más intenso es el CCB. Según \citet{schemm2014linkage}, el CCB se desplaza a baja altura a lo largo del lado frío del frente curvado (\textit{bent-back front}), experimentando un aumento de vorticidad potencial al pasar por debajo de la región de máxima liberación de calor latente en niveles medios —generada por el WCB ascendente—, y llega con alta vorticidad potencial al extremo del frente curvado, donde se localiza el jet de bajo nivel más intenso. Esto establece un acoplamiento indirecto y coherente entre el WCB (dominante en la distribución de precipitación) y el CCB (dominante en la distribución de vientos máximos), explicando por qué ambas variables presentan firmas espaciales asimétricas pero dinámicamente relacionadas. En el Atlántico Sur, \citet{bitencourt2010relating} y \citet{cardoso2022synoptic} documentan que los vientos máximos tienden a concentrarse en sectores específicos, como el cuadrante noreste. Además, las ráfagas más destructivas suelen estar vinculadas a características de mesoescala cuya posición relativa al centro de vorticidad varía a medida que el ciclón madura \citep{eisenstein2023identification}. Esta heterogeneidad explica por qué la intensidad de los máximos del viento no puede capturarse mediante una única métrica central \citep{corner2025classification}.

\end{itemize}

Esta disociación entre la métrica dinámica central y el impacto en superficie tiene consecuencias directas para el diseño analítico. \citet{andrade2024composite} documentan que, en el Atlántico Sudoccidental, los ciclones más intensos según MSLP presentan trayectorias sistemáticamente desplazadas hacia el sur respecto a los más intensos según viento máximo a 10 m, y que solo el 56\% de los sistemas que superan ambos umbrales (p10 de presión y p90 de viento) simultáneamente afectan la zona costera. Este resultado evidencia que el análisis de los extremos de viento requiere un marco de referencia centrado en el sistema —como el que proporciona el PDF bidimensional— y no un enfoque euleriano fijo, ya que la asimetría espacial de los máximos varía entre sistemas de igual intensidad dinámica central.

Esta disociación entre la métrica dinámica central y el impacto en superficie tiene consecuencias directas para el diseño analítico. \citet{andrade2024composite} documentan que, en el Atlántico Sudoccidental, los ciclones explosivos tienden a presentar trayectorias menos zonales a medida que aumenta su intensidad, con desplazamientos preferenciales de noroeste a sudeste e impactos  en las zonas costeras del sur de Brasil y Uruguay. Este resultado evidencia que el análisis de los extremos requiere un marco de referencia centrado en el sistema que permita extraer características generales de su estructura. Bajo esta premisa, autores como \citet{priestley2022improved} y \citet{han2025system} proponen el uso de la Estimación de Densidad de Kernel (KDE), el uso de esta herramienta permite caracterizar la asimetría espacial de los máximos y asi capturar la distribución sin perder información espacial de los ciclones.

\begin{equation}
\hat{f}(x) = \frac{1}{nh} \sum_{i=1}^{n} K\left(\frac{x - X_i}{h}\right)
\end{equation}

Esta técnica transforma observaciones discretas en una superficie de probabilidad continua, permitiendo mapear la \textit{huella} física del ciclón. Su aplicación es especialmente útil al segmentar el análisis por etapas evolutivas, aprovechando asi los resultados del algoritmo \textit{CycloPhaser} \citep{coutodesouza2024new}, ya que esto puede revelar cómo los campos extremos cambian de ubicacion o extensión a lo largo del ciclo de vida del sistema.

\subsection{Extremos de viento y precipitación}\label{subsec:tb_compound_def}

Históricamente, la intensidad de los ciclones se ha caracterizado mediante métricas puntuales, como la presión mínima central o la vorticidad máxima. 
Sin embargo, los efectos de viento y precipitación pueden actuar como fenómenos meteorológicos simultáneos; es aquí donde radica la relevancia de su concurrencia. 
Según argumentan \citet{zscheischler2018future}, los procesos que originan eventos intesos frecuentemente interactúan y presentan dependencias espaciales y temporales, lo que implica que su ocurrencia conjunta es físicamente dictada por el sistema al adoptar el enfoque de \textit{compound events} (eventos compuestos). En el contexto de la dinámica ciclónica extratropical, el caso de esta tesis corresponde a la co-ocurrencia espacial y temporal de viento y precipitación extremos dentro de una misma fase, cuya dependencia estadística —lejos de ser aleatoria— está físicamente dictada por la arquitectura tridimensional del ciclón (WCB y CCB).

En el contexto de la dinámica extratropical, \citet{chen2025characteristics} establecen que la interacción entre el campo de viento y la precipitación no es lineal ni espacialmente homogénea. 
La ocurrencia conjunta de estos extremos responde a mecanismos físicos acoplados: la intensificación dinámica suele estar modulada por procesos diabáticos (liberación de calor latente), lo que sugiere una dependencia estadística entre ambas variables que varía según la fase de vida del sistema. 

No obstante, es necesario reconocer una limitación en la cuantificación de estas variables. 
Autores como \citet{chen2024evaluation} realizaron evaluaciones de ERA5 sobre variables de superficie asociadas a ciclones extratropicales en el Hemisferio Norte, y demuestran que, si bien el producto presenta un sesgo general reducido para el viento (-0.7\%) y la precipitación (-10.4\%), este sesgo se amplifica para los eventos más intensos: para el 5\% de ETCs más extremos, la subestimación del viento alcanza el -10.2\% y la de la precipitación el -22.6\%. 
Esta tendencia a suavizar los máximos más intensos implica que los resultados obtenidos para el subconjunto p90 deben interpretarse como estimaciones conservadoras de la severidad de los sistemas.

\subsection{Descomposición modal de campos (EOF/MEOF)}\label{subsec:tb_math_meof}

Para aislar patrones coherentes dentro de este acoplamiento físico, el análisis de Funciones Ortogonales Empíricas (EOF) y su extensión multivariada (MEOF) se han establecido como el estándar metodológico. Según \citet{hannachi2007empirical}, esta técnica permite descomponer la variabilidad de un campo en modos ortogonales de varianza máxima, separando la señal física robusta del ruido de mesoescala. 

Específicamente para eventos compuestos, \citet{liang2018multivariate} demostraron la eficacia del MEOF al aplicarlo sobre un vector de estado normalizado $\mathbf{Z}$, que permite integrar métricas con unidades físicas dispares (como la precipitación y el viento superficial). Esta normalización estadística —cuya formulación matemática detallada y aplicación específica para este estudio se desarrolla exhaustivamente en la Sección \ref{subsec:eof_multivariado} de la Metodología— permite identificar si los máximos de ambas variables covarían en el mismo espacio geométrico o si, por el contrario, presentan desfases estructurales sistemáticos.

\begin{equation}
\mathbf{Z}(t) = \left[ \frac{X'*{prep}(t)}{\sigma*{prep}}, \frac{X'*{viento}(t)}{\sigma*{viento}} \right]
\end{equation}

Esta normalización permite identificar si los máximos de precipitación y viento covarían espacialmente o si presentan desfases estructurales sistemáticos. Físicamente, este grado de acoplamiento está dictado por la arquitectura tridimensional del ciclón; como evidencian \citet{catto2015front} y \citet{schemm2014ccb}, la precipitación extrema se concentra primordialmente en la zona de ascenso del WCB, mientras que los extremos de viento en superficie son modulados por la evolución del sector frío y el descenso del CCB. 
Este diseño sigue la práctica establecida por estudios de estructura ciclónica como el de \citet{priestley2022improved}, quienes aplican compositing centrado en el sistema al 10\% de ciclones más intensos (en términos de vorticidad) para aislar los patrones estructurales representativos.
